\section{Discussion}

\subsection{Cross-Validation Performance}

SVD achieved the best result on the dataset with an average RMSE of 0.4872, followed closely by the neural network at 0.8672. The matrix factorization and kNN
algorithms produced intermediate results at 0.9413 and 0.9645 respectively, and the random baseline established a baseline RMSE of 1.4289. The neural network's additional
capacity did not provide an advantage of the classical methods, likely due the size of the dataset and in turn the overfitting constraints that the network had to contend with,
as well as the fact that this dataset is known to be well-suited to linear factor models. Nonetheless, the matrix factorization methord performed significantly worse than SVD,
indicating that the bias terms and global mean setting of the predictions are valuable for modeling the user-item interractions within the data. The kNN method performed worse than
the other collaborative filtering methods, which was expected given how limiting the tag-based similarity is relative to the complexity of the rating signal.

\subsection{Pairwise RMSE}

As expected, the largest pairwise RMSE values involve the random algorithm, which given the performance measures indicates that the prediction behavior is quite different from the 
trained algorithms. Among the other four, SVD and neural network produce the most similar predictions at 0.2347, which suggest that the neural network is learning a similar latent factor 
representation of the rating data; however, given its performance, it is likely not as well-optimized due to its respective constraints. The SVD and matrix factorization algorithms
are noticeably different at 0.4348, though the absence of the bias terms produces a noticeably degraded and different prediction signal. The two most different algorithms aside from the baseline
were the kNN and matrix factorization at 0.6537, which utilize completely different signals and assumptions in producing a rating model.